\section{Résultats}

\subsection{Oscillations libres}

Avant de mesurer une oscillation, nous avons mesure le rapport de force qu'il faut exercer sur le balancier afin de le placer a un certain angle.
\\
\noindent \begin{minipage}{0.5\textwidth}
\includegraphics[width=\textwidth]{raideur.png}
	\captionof{figure}{$M_r$($\theta$)}
\label{fig:C}
\end{minipage}
\begin{minipage}{0.4\textwidth}
	Nous obtenons par consequent le graphique (\ref{fig:C}) ci-joint.
	Nous avons pu directement extraire la valeur C, soit la constante de raideur du ressort utilise grace a la relation (\ref{eq:nature}).\\ Soit :
	\\      C = 0.0123 \pm 0.0002 [Nm / rad]
	
\noindent \end{minipage}
\\
\\
\\
En un second temps, apres avoir montee l'installation pour mesures les oscillations du balencier-spiral (voir fig:\ref{fig:samplesetup}), nous avons donne une position hors-equilibre au balencier et mesurer en fonction du frein imposee a les differences oscillations observable en annexe (voir fig:\ref{fig:oscillations}).

\begin{minipage}[t]{0.45\textwidth}
    \centering
    \includegraphics[width=\linewidth]{lambda_quadratique.png}
	\captionof{figure}{$\lambda$(I)}
	\label{fig:lambda-current_quad}
\end{minipage}
\hfill
\begin{minipage}[t]{0.5\textwidth}
    \centering
    \includegraphics[width=\linewidth]{lambda.png}
	\captionof{figure}{$\lambda$($I^2$)}
	\label{fig:lambda-current}
\end{minipage}
\\
\\
La relation entre le facteur d'amortissement et le courant circulant dans les freins n'est pas lineaire. Nous constatons verifions la relation (\ref{eq:lambda-courant}), avec le graphique \ref{fig:lambda-current}. Nous pouvons extraire grace a la relation Eq. \ref{eq:w} la moyenne des pulsations propres $\overline{\omega_0}$ = 3.25 +/- 0.02 [rad/s] desquelles on extrait la frequence propre : $\overline{f_0} = 0.517 \pm 0.003 [Hz]$ partir de l'inverse de la periode : Eq. \ref{eq:t} (voir valeurs : \ref{val:f0}). Grace a la relation Eq~.\ref{eq:J} nous trouvons le moment d'inertie du balencier $J$ = $C$ / $\omega_^2$ = $0.0123/3.25^2$ = (1.165 \pm 0.024) \cdot $10^-3$ [Kg/$m^2$]. 

\newpage
\subsection{Expérience 2}

En tracant la phase en fonction de la frequence (voir \ref{fig:phase}), nous pouvons retrouver les frequences propres $f_0$ pour les alimentations a 0.5 en bleu et 0.9A en rouge.
        Les frequences propres resultent du point d'inflexions de leurs courbes correspondantes.
\\
\noindent \begin{minipage}{0.65\textwidth}
\vspace{0.4cm}
\includegraphics[width=\textwidth]{force_phase.png}
	\captionof{figure}{$\phi$(f), evolution de la Phase}
\label{fig:phase}
\end{minipage}
\hspace{0.00005\textwidth}
\begin{minipage}{0.4\textwidth}
	Ci-joint (\ref{fig:phase}), les frequences propres ont pu s'obtenir a partir de la relation de la phase-freqence (voir \ref{eq:phase_forcee}), sachant que la phase en fonction de la pulsation propore vaut $\pi/2$. L'ajustement nous donne :
	\\ \\ $f_0$(0.5A) = 0.518 \pm 0.009 [Hz].
	\\ \\ $f_0$(0.9A) = 0.528 \pm 0.004 [Hz].
\noindent \end{minipage}

\noindent \begin{minipage}{0.65\textwidth}
\vspace{0.4cm}
\includegraphics[width=\textwidth]{force_amp.png}
	\captionof{figure}{A(f), evolution de l'elongation angulaire}
	\label{fig:amplitude}
\label{fig:nature}
\end{minipage}
\hspace{0.00005\textwidth}
\begin{minipage}{0.4\textwidth}
	En revanche, l'amplitude du balencier en fonction de la frequence nous permet d'obtenir les frequences propres en fonction de Eq. \ref{eq:omega_res}, soit en explorant le lien entre pulsation propre et pseudo pulsation :
	\\ \\ $f_0$(0.5A) = 0.518 \pm 0.004 [Hz].
	\\ \\ $f_0$(0.9A) = 0.519 \pm 0.001 [Hz]
	
\noindent \end{minipage}
\\

En fonction du graph d'amplitude,
nous obtenons en realisant un fit a partir du modele Eq. \ref{eq:amplitude_forcee} les valeurs suivantes :
\\\\
\textbf{$\lambda$(0.5A) = 0.191 ± 0.003 [$1/s$] , \hspace{2.5cm} $\lambda$(0.9A) = 0.609 ± 0.005 [$1/s$]. \\$\omega_0$(0.5A) = 3.249 ± 0.001 [rad/s],  \hspace{2cm} $\omega_0$(0.9A) = 3.315 ± 0.003 [rad/s]} \\\\Et la par moyenne des freqences obtenues : 
\\- \hspace{5cm} \textbf{$\overline{f_0}$ = 0.520 ± 0.005 [Hz].} 

\vspace{0.5cm}
A partir de la puissance relative, si nous etablissons donc le rapport $\overline{P}$ / $P_0$ en fonction de la frequence sous differents courants, nous pouvons observer une courbe plus aplatie (voir \ref{fig:power9}) qu'en le graphique a 0.5 A (voir \ref{fig:power5}).
Ces courbes vont nous permettre d'extraire FMWH grace a des outils MATLAB et les facteurs de qualite Q a partir de la relation  Eq. \ref{eq:Q_fwhm}.
\\
\begin{minipage}[t]{0.5\textwidth}
    \centering
	\includegraphics[scale=0.3]{power_0.5A.png}
	\captionof{figure}{$P$/$P_0$(f) 0.5A}
	\label{fig:power5}
\end{minipage}
\hfill
\begin{minipage}[t]{0.5\textwidth}
    \centering
	\includegraphics[scale=0.3]{power_0.9A.png}
	\captionof{figure}{$P$/$P_0$(f) a 0.9A}
	\label{fig:power9}
\end{minipage}

\vspace{0.5cm}
À partir du rapport de puissance, nous obtenons a partir du fit :

\[
	\mathrm{f_0}(0.5A) = 0.521 \pm 0.005~\mathrm{Hz}, \qquad
\mathrm{f_0}(0.9A) = 0.47 \pm 0.01~\mathrm{Hz}.
\]

Puis, nous obtenons la FWHM, directement liée à la largeur de la résonance :

\[
\mathrm{FWHM}_{0.5\,A} = 0.132 \pm 0.003~\mathrm{Hz}, \qquad
\mathrm{FWHM}_{0.9\,A} = 0.214 \pm 0.003~\mathrm{Hz}.
\]

Ce qui nous donne respectivement :
\[
	\mathrm{Q}_{0.5\A} = 0.521 / 0.132 = 4.0 \pm 0.1 \qquad
	\mathrm{Q}_{0.9\A} = 0.47 / 0.214 = 2,20 \pm 0.06
\]


%%%%%%%%%%%%%%%%%%%%%%%%%%%%%%%
%	ANALYSE
%%%%%%%%%%%%%%%%%%%%%%%%%%%%%%%

\section{Analyse des résultats, complement}

\subsection{oscillations libres}

Nous avons vu pour la constante de torsion $C = 0.0123 \pm 0.0002 [Nm / rad]$ que l'incertitude est relativement basse. Les valeurs sont obtenu directement grace a la pente du la fonction $M_r$($\theta$) (voir fig: \ref{fig:C}).

\\

En revanche, si l'on souhaite déterminer l'incertitude sur le moment d'inertie du balancier,
nous pouvons calculer l'erreur sur $J$ à partir de la relation entre $C$, $J$ et la pulsation propre
$\omega_0$ (voir Eq.~\ref{eq:J}), en utilisant la méthode des dérivées partielles avec :

\[
\frac{\partial J}{\partial C} = \frac{1}{\omega_0^2},
\qquad
\frac{\partial J}{\partial \omega_0}
= -\,\frac{2C}{\omega_0^3}.
\]

L'incertitude sur $J$ s'écrit donc :

\[
\sigma_J = 
\sqrt{
\left( \frac{\sigma_C}{\omega_0^2} \right)^2
+
\left( \frac{2C\,\sigma_{\omega_0}}{\omega_0^3} \right)^2
}.
\]

En remplaçant par les valeurs numériques utilisées :

\[
\sigma_J =
\sqrt{
\left( \frac{0.0002}{3.25^2} \right)^2 +
\left( \frac{2 \cdot 0.0123 \cdot 0.02}{3.25^3} \right)^2
}
= 0.0237 \times 10^{-3}\ \mathrm{kg\,m^2}
\]


\subsection{oscillations forcees}

Le facteur de qualité est donné par (Eq.~\ref{eq:Q}) et l’incertitude correspondante, obtenue par dérivées partielles, est :

\[
\frac{\partial Q}{\partial f_0} = \frac{1}{\mathrm{FWHM}},
\qquad
\frac{\partial Q}{\partial\,\mathrm{FWHM}}
= -\,\frac{f_0}{\mathrm{FWHM}^2}.
\]

Ainsi,

\[
\sigma_Q
=
\sqrt{
\left(\frac{\sigma_{f_0}}{\mathrm{FWHM}}\right)^2
+
\left(
\frac{f_0\,\sigma_{\mathrm{FWHM}}}{\mathrm{FWHM}^2}
\right)^2
}.
\]

\[
\sigma_{Q_{0.5A}}
=
\sqrt{
\left(\frac{0.005}{0.132}\right)^2
+
\left(\frac{0.521 \cdot 0.003}{0.132^2}\right)^2
}
= 0.097
\]

\[
\sigma_{Q_{0.9A}}
=
\sqrt{
\left(\frac{0.01}{0.214}\right)^2
+
\left(\frac{0.47 \cdot 0.003}{0.214^2}\right)^2
}
= 0.056
\]

%=====================================
%%      CONCLUSIONS
%=====================================

\section{Conclusion}

%La constante de torsion C est de $0.0123 \pm 0.0002 [Nm / rad]$.     % OK
%\\Le moment d'inertie J est de $(1.165 \pm 0.024)$ \cdot $10^-3$ [Kg/$m^2$] % OK
%\\La moyenne des frequence de resonance $\overline{$f_0$}$ du systeme en regime libre est de $0.517 \pm 0.003$ % OK
%\\La frequence de resonance du systeme en dephasage (methode des phaseurs) est de $0.520 \pm 0.005 [Hz]$ % OK
%\\La relation $\lambda$($I^2$) est de $0.7801 \pm 0.0026$ [$s^-1$$A^-2$] 
%\\Les facteurs de qualite des systemes sont :
%\\      $Q_0.5A$ =  4.0 \pm 0.1 
%\\      $Q_0.9A$ =  2.20 \pm 0.06 

L’étude menée sur le pendule de Pohl nous a permis de caractériser avec précision les paramètres essentiels du système. 
La constante de torsion a été évaluée à 
\[
C = 0.0123 \pm 0.0002~\mathrm{N\,m/rad},
\]
ce qui a conduit à une détermination fiable du moment d’inertie du balancier :
\[
J = (1.165 \pm 0.024)\times 10^{-3}~\mathrm{kg\,m^2}.
\]

L’analyse des oscillations libres a permis d’obtenir une fréquence propre moyenne de
\[
\overline{f_0} = 0.517 \pm 0.003~\mathrm{Hz},
\]
tandis que la méthode des phaseurs a fourni une valeur tout à fait cohérente :
\[
f_0 = 0.520 \pm 0.005~\mathrm{Hz}.
\]

L’évolution de l’amortissement en fonction du courant s’est révélée conforme au modèle attendu, avec la relation
\[
\lambda(I^2) = 0.7801 \pm 0.0026~\mathrm{s^{-1}A^{-2}}.
\]

Les facteurs de qualité mesurés pour deux régimes d’excitation illustrent clairement la diminution de la résonance lorsque l’amortissement augmente :
\[
Q_{0.5\,\mathrm{A}} = 4.0 \pm 0.1, \qquad
Q_{0.9\,\mathrm{A}} = 2.20 \pm 0.06.
\]

Comme prévu pour un oscillateur harmonique forcé, l’amplitude atteint son maximum à la résonance et le déphasage s’établit alors à $-90^\circ$. 
Notons enfin que la précision des ajustements est restée limitée par le manque de points de mesure autour du pic de résonance, ce qui laisse entrevoir des marges d’amélioration dans l’acquisition de données futures.
%=====================================
%       SOURCES
%=====================================

%\begin{thebibliography}{99}

%\bibitem{schemapohl}
%Gilbert Gastebois, \textit{Le pendule de Pohl},
%\\
%                \texttt{https://ggastebois.fr/java/pendpohl/theorie\_pend\_pohl.pdf}
%\\
%        Consulte le : 17 Novembre 2025

%\end{thebibliography}
